\section{括区与二分法}
\label{sec:08-Numerical-Analysis/02-Root-Finding/01}

你在标定一个定价模型的折现率。模型是一段跑二十秒的仿真代码，输入一个折现率，输出一个现金流缺口。要找的就是让缺口归零的那个输入。没有解析式，导数拿不到，代码在某些极端参数上还会直接抛异常。

手上只有两个数：折现率取 \(1\%\) 时缺口为正，取 \(8\%\) 时为负。

听上去太单薄了。可只要缺口随折现率连续变化，这两个数已经足以启动一次必然成功的搜索——而且在第一次调用仿真之前，你就能说出最多要跑多少次。

\subsection{一次变号就是一张含根证书}

\begin{proposition}
若 \(f\) 在闭区间 \([a,b]\) 上连续，且 \(f(a)f(b)<0\)，则存在 \(r\in(a,b)\) 使得 \(f(r)=0\)。
\end{proposition}

这是介值定理的直接推论，来历见\hyperref[sec:02-Calculus/02-Limits-and-Continuity/04]{``连续性与整体保证''}。两个条件各自防住一种失效，值得分开对账。

去掉连续性，符号可以靠一次跳跃翻转而不经过零：\(f(x)=-1\)（\(x<0\)）与 \(f(x)=1\)（\(x\ge 0\)）在 \([-1,1]\) 上端点痛快地异号，函数却从不取零。去掉变号，函数可以擦着横轴折返而端点同号：\(f(x)=(x-1)^2\) 在 \(x=1\) 处触地弹回，包含 \(1\) 的任何区间两端都是正的。连续性把``符号变了''翻译成``穿过了横轴''，变号则保证确实有得穿；缺哪一个，推理链都断在原地。

有了这张证书，算法本身几乎没有内容：取中点 \(m=(a+b)/2\)，算 \(f(m)\)，把 \(m\) 与两端中符号相反的那个配成新区间，另一半丢掉。函数值到底是 \(10^{-9}\) 还是 \(10^{7}\) 无关紧要，只有正负号被使用。每一步之后，``根在当前区间内''这句话仍然成立——它是一个不变量，从头到尾没有被交易出去。

\begin{center}
\begin{tikzpicture}[x=1cm,y=1cm,font=\small,>=stealth]
  \draw[->,black!55] (-0.35,0) -- (7.9,0) node[right,font=\footnotesize] {\(x\)};
  \draw[blue!75,thick] plot[domain=1:2,samples=60] ({7*(\x-1)},{0.45*(\x*\x*\x-\x-2)});
  \draw[black!30,dashed] (3.65,2.1) -- (3.65,-4.55);
  \node[above,font=\footnotesize,black!60] at (3.65,2.05) {真根};
  \fill[black!70] (0,-0.9) circle (1.4pt);
  \fill[black!70] (7,1.8) circle (1.4pt);
  \node[left,font=\footnotesize] at (-0.05,-0.9) {\(f(a)<0\)};
  \node[right,font=\footnotesize] at (7.05,1.75) {\(f(b)>0\)};
  % 区间条
  \draw[black!65,line width=1.1pt] (0,-1.9) -- (7,-1.9);
  \draw[black!65] (0,-2.05) -- (0,-1.75);
  \draw[black!65] (7,-2.05) -- (7,-1.75);
  \draw[red!70] (3.5,-2.05) -- (3.5,-1.75);
  \node[right,font=\footnotesize,black!70] at (7.25,-1.9) {\(L_0\)};
  \draw[black!65,line width=1.1pt] (3.5,-2.6) -- (7,-2.6);
  \draw[black!65] (3.5,-2.75) -- (3.5,-2.45);
  \draw[black!65] (7,-2.75) -- (7,-2.45);
  \draw[red!70] (5.25,-2.75) -- (5.25,-2.45);
  \node[right,font=\footnotesize,black!70] at (7.25,-2.6) {\(L_0/2\)};
  \draw[black!65,line width=1.1pt] (3.5,-3.3) -- (5.25,-3.3);
  \draw[black!65] (3.5,-3.45) -- (3.5,-3.15);
  \draw[black!65] (5.25,-3.45) -- (5.25,-3.15);
  \draw[red!70] (4.375,-3.45) -- (4.375,-3.15);
  \node[right,font=\footnotesize,black!70] at (7.25,-3.3) {\(L_0/4\)};
  \draw[black!65,line width=1.1pt] (3.5,-4.0) -- (4.375,-4.0);
  \draw[black!65] (3.5,-4.15) -- (3.5,-3.85);
  \draw[black!65] (4.375,-4.15) -- (4.375,-3.85);
  \draw[red!70] (3.9375,-4.15) -- (3.9375,-3.85);
  \node[right,font=\footnotesize,black!70] at (7.25,-4.0) {\(L_0/8\)};
  \draw[black!65,line width=1.1pt] (3.5,-4.7) -- (3.9375,-4.7);
  \draw[black!65] (3.5,-4.85) -- (3.5,-4.55);
  \draw[black!65] (3.9375,-4.85) -- (3.9375,-4.55);
  \node[right,font=\footnotesize,black!70] at (7.25,-4.7) {\(L_0/16\)};
\end{tikzpicture}

\emph{曲线是 \(f(x)=x^3-x-2\) 在 \([1,2]\) 上的图像，虚线是它的根。下方每一条横杠是一次对折之后剩下的区间，红色短竖线标出该步测试的中点。四步之后区间只剩原来的十六分之一，而``根在杠内''这句话一次也没有失效。}
\end{center}

\subsection{误差界在第一次调用之前就算得出来}

记初始宽度 \(L_0=b_0-a_0\)。每步对折，\(n\) 步后 \(L_n=L_0/2^{\,n}\)。取当前区间的中点 \(m_n\) 作近似，根一定落在同一区间里，于是

\[
\abs{m_n-r}\le\frac{L_n}{2}=\frac{L_0}{2^{\,n+1}}.
\]

这个不等式里没有出现 \(f\) 的任何性质——没有导数，没有曲率，没有 Lipschitz 常数。它只用到了区间宽度。所以要把位置误差压到容差 \(\tau\) 以下，所需步数

\[
n\ \ge\ \left\lceil\log_2\frac{L_0}{\tau}\right\rceil
\]

在你按下运行键之前就是已知的。折现率那个例子里 \(L_0=7\%\)，要精确到 \(10^{-6}\)，\(\log_2(7\times10^{4})\approx 16.1\)，十七步封顶：仿真跑十七次，答案带一张误差证书。

代价写在同一个公式里。每步只买到一个二进制位，约 \(0.30\) 个十进制位。想要十六位有效数字，得跑五十多步。下一节的 Newton 法在同样的问题上通常五六步就撞上机器精度——那才是二分法真正的对手。

有一个实现细节容易被跳过。中点直接写成 \((a+b)/2\)，在 \(a,b\) 同号且接近浮点上界时会溢出；写成 \(a+(b-a)/2\) 稳妥得多。另一个信号是：算出的中点已经等于某个端点。这说明当前精度下区间不能再切了，继续迭代只会空转，此时应当停下并报告已达精度极限，而不是等一个永远不会满足的容差。同理，判断异号最好分别取符号而不是算乘积 \(f(a)f(b)\)——两个大数相乘会上溢，两个小数相乘会下溢到零，而符号本身从不溢出。

\subsection{变号是充分条件，不是必要条件}

上面那个 \(f(x)=(x-1)^2\) 不只是个反例，它代表一整类情形：偶重根两侧函数值同号，二分法对它完全失明。方法是靠符号变化找根的，这里根本没有符号变化，再密的初始搜索也扫不出括区。奇重根没有这个问题——三重根 \((x-1)^3\) 照样穿轴。

区间内藏了多个根也一样。二分法会收敛到其中某一个，但不会告诉你还有几个同伴留在被丢掉的那半边。它不枚举，只收敛。要枚举，得先做粗扫描把根分离到各自的括区里，这本身是另一个问题。

这两条限制放在一起，给出了二分法的适用边界：它处理的是``已知有一次穿轴，把它定位到给定精度''，不是``找出所有根''，也不是``判断有没有根''。

\subsection{残差小不等于离根近}

一个常见的收工判据是 \(\abs{f(x)}<\varepsilon\)。它比看上去脆弱。在根附近做一阶展开，

\[
\abs{x-r}\ \approx\ \frac{\abs{f(x)}}{\abs{f'(r)}},
\]

位置误差是残差除以斜率。斜率很小时，一个离根还很远的点也能交出漂亮的小残差；斜率很大时，一个位置极准的点可能残差不小，被误判成没收敛而白白丢掉。

\begin{center}
\begin{tikzpicture}[x=1cm,y=1cm,font=\small,>=stealth]
  \fill[black!8] (-3.3,-0.34) rectangle (3.3,0.34);
  \draw[->,black!55] (-3.5,0) -- (3.7,0) node[right,font=\footnotesize] {\(x\)};
  \draw[black!55] (0,-1.15) -- (0,1.25);
  \node[above,font=\footnotesize,black!60] at (0,1.25) {根};
  \draw[blue!75,thick] (-3.2,-0.70) -- (3.2,0.70);
  \node[right,font=\footnotesize,blue!75] at (3.24,0.70) {平缓};
  \draw[red!75,thick] (-0.6,-0.96) -- (0.6,0.96);
  \node[right,font=\footnotesize,red!75] at (0.62,0.99) {陡峭};
  \node[right,font=\footnotesize,black!65] at (-3.45,0.62) {\(\abs{f}\le\varepsilon\)};
  \draw[black!35,dotted] (-1.55,-0.34) -- (-1.55,-1.4);
  \draw[black!35,dotted] (1.55,-0.34) -- (1.55,-1.4);
  \draw[blue!70,line width=1.2pt] (-1.55,-1.4) -- (1.55,-1.4);
  \draw[black!35,dotted] (-0.21,-0.34) -- (-0.21,-1.85);
  \draw[black!35,dotted] (0.21,-0.34) -- (0.21,-1.85);
  \draw[red!70,line width=1.2pt] (-0.21,-1.85) -- (0.21,-1.85);
  \node[right,font=\footnotesize,black!70] at (1.75,-1.4) {残差达标的 \(x\) 范围（平缓）};
  \node[right,font=\footnotesize,black!70] at (1.75,-1.85) {同一容差下（陡峭）};
\end{tikzpicture}

\emph{同一个根，同一条残差容差带 \(\abs{f}\le\varepsilon\)。平缓函数被这条带截出一段很宽的横轴区间，陡峭函数只截出很窄的一段。残差达标说明的不是``离根多近''，而是``离根多近乘以斜率''。}
\end{center}

把这张图和上一张对照着看：二分法的误差界画在横轴上，残差判据的误差界画在纵轴上，只有知道 \(\abs{f'(r)}\) 才能在两者之间换算。而 \(1/\abs{f'(r)}\) 正是求根问题的条件数，它衡量的是输入扰动被放大成根的位移的倍数，与\hyperref[sec:08-Numerical-Analysis/01-Floating-Point-and-Error/03]{``条件数衡量问题本身的敏感性''}里的量是同一个。二分法的可贵之处在于它绕开了这次换算：区间宽度直接就是位置误差。

实践中通常同时检查区间宽度是否达到绝对与相对容差、中点残差是否足够小、以及中点是否已经与端点重合。单独任何一条都不构成收敛的证据。

\subsection{什么时候值得选慢的那个}

二分法适合的场合远不止教科书里的多项式。单调仿真的反演——给定目标产出 \(y\)，找输入 \(x\) 使 \(f(x)=y\)——只需要函数值接口。分位数定位是同一件事，分布函数单调，二分天然匹配。事件过零的判定也是：某个监控量随时间从正变负，你要知道它在哪一刻穿过阈值，精度要求不高但绝不能漏报。

更常见的用法是当安全网。先用二分建立一个可靠的括区，再在区间内尝试快速步；候选点落在区间内且让区间显著缩小就接受，否则退回一次对折。这样最坏情况仍由对折兜底，平均情况由快速步推进。下一节和再下一节会把这套结构讲完整。

判断的落点其实很简单：当求根失败的代价高于慢一点的代价时，选二分。它嵌在另一个循环内部时尤其如此——外层优化不会告诉你内层求根悄悄发散了，你只会看到一堆莫名其妙的结果。

不过每步 \(0.3\) 个十进制位确实太贵了。二分法之所以慢，是因为它只用了 \(f\) 的符号，把函数值的大小、图像的走向、导数的信息全部扔掉。下一节换一种交易：假设 \(f\) 光滑、导数可得，用切线代替曲线，看看能把速度提到什么程度，又要为此交出哪些保证。
